%! Author = sriadityavedantam
%! Date = 3/29/26

\section{Sets and Classes}

Set theory is very much its own field, so we will not be getting into the specifics and the nitty-gritty of each topic.
It will just be a brief overview.

\begin{definition}[Sets and Basic Set-Theoretic Language]
    \textbf{Elements} are either a part of a set or not part of a set.
    There are infinitely many elements, and they have a choice of being a member of a set.
    When an element $x$ is a member of a set $A$, we denote this by
    \[
        x \in A.
    \]
    Otherwise, we say
    \[
        x \notin A.
    \]
    We can also write this out in words as ``$x$ is (not) an element of $A$.''
    These are some of the few things most people use symbols for regardless of their preferences for symbolic language.
    The following are predicates.

    \textbf{1. ``For all''}
    This is denoted by $\forall$.

    \textbf{2. ``There exists''}
    This is denoted by $\exists$.
    The \textbf{axiom of extensionality} states that given sets $A$ and $B$, if for all elements $x$ we have
    \[
        x \in A \iff x \in B.
    \]
    Then $A = B$.
    If for all elements $x$, whenever $x \in A$ we also have $x \in B$, then $A$ is a \textbf{subset} of $B$, denoted by $A \subseteq B$.

    The \textbf{empty set} is a set with no elements, denoted by $\varnothing$.

    A \textbf{class of sets} is a collection that contains sets and only sets.

    The \textbf{power axiom} states that for every set $A$, the power class $P(A)$ contains all subsets of $A$ within a set.
    This is often denoted by $2^A$ and has $2^{|A|}$ elements.

    A \textbf{union of sets} considers all the elements in both sets, denoted by $A \cup B$.

    An \textbf{intersection of sets} considers only the common elements in both sets, denoted by $A \cap B$.

    A \textbf{disjoint set} situation occurs when $A \cap B = \varnothing$.

    A \textbf{family of sets} is a class of sets where each element, mind you a set, is indexed.
    This is generally denoted by
    \[
        \bigcup_{i \in I} A_i \coloneqq \{x : x \in A_i \text{ for some } i \in I\}.
    \]
    Similarly, we define
    \[
        \bigcap_{i \in I} A_i.
    \]

    The \textbf{complement} of $A$ is related to the negation of $A$, where we use DeMorgan's Laws.
\end{definition}